
\def\PUDcodigo{ECA224}

\def\PUDcodacad{04507.41}

\def\PUDdisciplina{Acionamento de Máquinas}

\def\PUDsemestre{S8}

\def\PUDhoras{80}

%NOVOS ---------------------------------------------------
\def\PUDhorasTeoria{60}
\def\PUDhorasPratica{20}

\def\PUDinterECA{ 
\hyperref[PUD:ECA216]{Eletrônica III (04507.29)}

\hyperref[PUD:ECA222]{Instalações Elétricas (04507.30)}

\hyperref[PUD:ECA233]{Controle I (04507.31)}

\hyperref[PUD:ECA223]{Máquinas Elétricas (04507.36)} }

\def\PUDatuaisECA{- Aplicações de acionamentos de máquinas na geração de energia eólica}

\def\PUDrecursos{Projetor multimídia; Quadro branco e pincel; Aulas em laboratório de informática e de acionamentos de máquinas.}

%----------------------------------------------------------

\def\PUDcreditos{4}

\def\PUDprerequisito{\hyperref[PUD:ECA233]{Controle I (04507.31)}}

\def\PUDpreacad{\hyperref[PUD:ECA233]{Controle I (04507.31)}}

\def\PUDobjetivos{Conhecer o controle de máquinas CC, MCEs e máquinas CA. Conhecer os principais métodos para variação da velocidade de MITs.  Parametrizar e controlar o conjugado e a velocidade de MITs via inversores industriais. Parametrizar e controlar a velocidade e a posição de servomotores.  Controlar a velocidade de motores CC.
%Conhecer os principais métodos para redução da corrente de partida em MITs.  Conhecer os principais métodos para variação da velocidade de MITs.  Parametrização e controle de conjugado e velocidade de MITs via inversores industriais.  Parametrização e controle de velocidade e posição de servomotores.  Controle de velocidade de motores CC.
}

\def\PUDementa{Introdução aos Acionamentos Elétricos; Sistemas Mecânicos e Acionamentos Elétricos; Eletrônica de Potência em Acionamentos; Acionamentos de Motores CC, MCEs e Máquinas CA; Controle de Malha Fechada para Acionamentos; Métodos de Sintonia MO e SO; Máquinas CA e Vetores Espaciais; Transformações entre Referenciais e Controle Vetorial. Parametrização de inversores industriais e controle de velocidade e posição de servomotores.
%Dispositivos de redução de corrente de partida e partida suave de MITs;  Dispositivos de variação de velocidade de MITs;  Acionamentos eletrônicos de máquinas elétricas rotativas;  Inversores industrial com controle escalar e vetorial.
}

\def\PUDconteudo{ & Introdução aos Acionamentos Elétricos: História; Elementos de um Acionamento Elétrico; Crescimento dos Acionamentos Elétricos; Aplicações Típicas de Acionamentos Elétricos; Multidisciplinaridade.
%Circuitos eletromecânicos de partida de MITs. 
\\ 
 & Sistemas Mecânicos e Acionamentos Elétricos: Sistemas com Movimento Linear; Sistemas Rotativos; Atrito; Torção no Eixo; Analogia Elétrica; Mecanismos de Acoplamento; Tipos de Cargas; Operação em Quatro Quadrantes; Operação em Regime Permanente e Dinâmico.
 %Controladores de tensão CA. 
\\ 
 & Eletrônica de Potência em Acionamentos: Semicondutores de Potência; Visão Geral das UPPs; Conversores para Motores CC; Síntese de CA de Baixa Frequência; Inversores Trifásicos; Outras Técnicas de Modulação.
 %Controle PWM e SPWM. 
\\ 
 & Acionamentos de Motores CC e MCEs: Estrutura das Máquinas CC; Princípios de Operação das Máquinas CC; Circuito Equivalente da Máquina CC; Modos de Operação; Enfraquecimento de Campo; UPPs em Acionamentos de Motores CC; Acionamentos de MCEs.
 %Inversores Industriais.  
\\ 
 & Controle de Malha Fechada para Acionamentos: Objetivos de Controle; Estrutura de Controle em Cascata; Passos no Projeto do Controlador Realimentado; Análise de Pequenos Sinais; Exemplo de Projeto do Controlador; Função da Alimentação Direta (\textit{feed-forward}); Efeito dos Limites; Integração Antissaturação (\textit{anti-windup}).
 %Inversores trifásicos VSI e CSI. 
 \\
 & Métodos de Sintonia MO e SO: Conceitos Iniciais; Sintonia de PIs pelo Método da Magnitude Ótima; Sintonia de PIs pelo Método Simétrico Ótimo
 \\
 & Máquinas CA e Vetores Espaciais: Distribuição Senoidal do Enrolamento do Estator; Representação por Vetores Espaciais; Excitação Senoidal e Equilibrada
 \\
 & Máquinas Síncronas de Ímãs Permanentes: Estrutura das Máquinas CA de Ímãs Permanentes; Princípio de Operação; Controlador e UPP
 \\
 & Transformações entre Referenciais: Transformada de Clarke; Transformada de Park; Transformações com Invariância em Amplitude; Transformações com Invariância em Potência.
 \\
 & Modelagem e Controle da MSIP: Considerações Iniciais; Modelo em Coordenadas abc; Modelo em Coordenadas dq; Controle Vetorial da MSIP.
 \\
 & Parametrização de inversores industriais e controle de velocidade e posição de servomotores.
 }

\def\PUDmetodo{Aulas expositórias que podem ser teóricas e/ou práticas, onde as práticas no laboratório serão marcadas no decorrer da disciplina. Além disso, são elaboradas simulações dos diversos acionamentos apresentados ao longo da disciplina.}

\def\PUDavalia{A avaliação será feita com aplicação de provas teóricas e/ou práticas, além da possibilidade de inclusão de trabalhos e seminários no decorrer da disciplina.}

\def\PUDbibbasica{ &  \href{www.biblioteca.ifce.edu.br/index.asp?codigo_sophia=40390}{Stephan}, Richard M.  Acionamento, Comando e controle de máquinas Elétricas.  Rio de Janeiro, RJ: Ciência Moderna, 2013.  230p.  ISBN: 9788539903542.
\\ 
&  \href{www.biblioteca.ifce.edu.br/index.asp?codigo_sophia=23846}{Franchi}, Claiton Moro.  Acionamentos Elétricos.  4º ed.  São Paulo, SP : Érica, 2011.  250p.  ISBN: 9788536501499.
\\ 
 &  \href{www.biblioteca.ifce.edu.br/index.asp?codigo_sophia=62576}{Maya}, Paulo A.; Leonardi, Fabrizio.  Controle essencial.  2º ed.  São Paulo, SP: Pearson Education do Brasil, 2014.  347p.  ISBN: 9788543002415.
 }

\def\PUDbibcomplementar{ & \href{www.biblioteca.ifce.edu.br/index.asp?codigo_sophia=23684}{Fitzgerald}, A. E; Kingley Jr, Charles; Umans, Stephen D.  Máquinas elétricas: com introdução à eletrônica de potência.  6º ed.  Porto Alegre, RS: Bookman, 2008.  648p.  ISBN: 9788560031047.
\\
 & \href{www.biblioteca.ifce.edu.br/index.asp?codigo_sophia=23744}{Leludak}, Jorge Assade.  Acionamentos eletromagnéticos.  Curitiba, PR: Base Editorial, 2010.  176p.  ISBN: 9788579055591.
\\
&  \href{www.biblioteca.ifce.edu.br/index.asp?codigo_sophia=62367}{Bim}, Edson.  Máquinas Elétricas e Acionamento.  3º ed.  Rio de Janeiro, RJ: Campus : Elsevier, 2014.  571p.  ISBN: 9788535259230.
\\
 & \href{www.biblioteca.ifce.edu.br/index.asp?codigo_sophia=65036}{Filippo Filho}, Guilherme.  Motor de indução: princípios de funcionamento, características operacionais, aplicações, acionamentos e comandos.  2º ed.  São Paulo, SP: Érica, 2016.  294p.  ISBN: 9788536504483.
\\
 & \href{www.biblioteca.ifce.edu.br/index.asp?codigo_sophia=40033}{Simone}, Gilio Aluisio.  Conversão eletromecânica de energia: uma introdução ao estudo.  São Paulo, SP: Érica, 2014.  324p.  ISBN: 9788571946033.
%\\
 %& \href{www.biblioteca.ifce.edu.br/index.asp?codigo_sophia=44617}{Rezek}, Ângelo José Junqueira.  Fundamentos básicos de máquinas elétricas: teoria e ensaios.  Rio de Janeiro, RJ: Synergia, 2011.  123p.  ISBN: 9788561325695.
%\\ 
 %&   Mohan, N;  Underland, T;   e Riobbins, W.  Power Electronics - converters, applications and design.  3 ed.  Editora Wiley \& Sons.  2006.  ISBN 9780471452058
%\\ 
 %&  IRWIN, J.  David.  Electrical Machines and Power Electronics.  Editora Taylor \& Francis USA.  2011.  350p.  ISBN : 9781439802854.
 }

\def\PUDautor{Adriano Holanda Pereira}

\def\PUDdata{2013-04-29}

\def\PUDrevisao{3}

\def\PUDrevisor{Celso Rogério Schmidlin Júnior}

\def\PUDdatarevisao{2019-05-15}

\PUDcorpo
